% Shared preamble: everything both the ICLR and the arXiv build need.
% Venue-specific bits (documentclass, style file, title block, bib style)
% live in build/iclr2027.tex and build/arxiv.tex.

\usepackage[utf8]{inputenc}
\usepackage[T1]{fontenc}
\usepackage{amsmath,amssymb,amsthm}
\usepackage{enumitem}
\usepackage{graphicx}
\usepackage{booktabs}
\usepackage{multirow}
\usepackage[table]{xcolor}
\usepackage{url}
\usepackage{hyperref}

% Verbatim inclusion of the pre-registration documents in the appendix.
% They are plain text and must appear exactly as committed, so they are
% \lstinputlisting'd from prereg/ rather than retyped: the appendix cannot
% drift from the files, because it is the files. Widest line in the set is
% 85 columns, which fits at \footnotesize, but breaklines is on so that a
% future edit to one of them cannot silently produce an overfull line.
% The section sign is the only non-ASCII character any of them contains.
\usepackage{listings}
% upquote before any listing: without it a backtick and an apostrophe come
% out as curly quotes in the tt font, which misrepresents documents that use
% backticks for code spans throughout.
\usepackage{textcomp}
\usepackage{upquote}
\lstset{
  basicstyle=\footnotesize\ttfamily,
  breaklines=true,
  breakindent=0pt,
  columns=fullflexible,
  keepspaces=true,
  showstringspaces=false,
  frame=none,
  xleftmargin=0pt,
  aboveskip=0.6\baselineskip,
  belowskip=0.6\baselineskip,
  % Only the section sign needs mapping. Do NOT add a {\S} key here: in a
  % literate list that is read as the single character S, and every capital
  % S in the documents comes out as a section sign.
  literate={§}{{\S}}1,
}
% Not microtype: under the XeTeX engine tectonic uses, only protrusion is
% available (no font expansion for these Type 1 faces), and enabling it
% moved the arXiv build from six overfull lines to eight. Measured, not
% assumed; if the build ever moves to pdflatex, retest it.

% article.cls ships \belowcaptionskip = 0pt, which puts the first \toprule on
% the baseline of the caption's last line. It only shows up on some tables,
% depending on how the last caption line falls, and it looks like a printing
% error when it does. Give every caption a little air instead.
\setlength{\belowcaptionskip}{6pt}

% Heatmap-style cell shading for the matrix tables.
% \best = winner of the row (lowest NLL excess / highest accuracy).
% \worst = loser of the row.
% NOTE: these need colortbl, i.e. xcolor's [table] option. Each build file
% issues \PassOptionsToPackage{table}{xcolor} on its FIRST line, because the
% ICLR style file pulls in eso-pic -> xcolor and would otherwise option-clash,
% silently leaving \cellcolor undefined and printing a literal "green!18".
\newcommand{\best}[1]{\cellcolor{green!18}\textbf{#1}}
\newcommand{\worst}[1]{\cellcolor{red!10}#1}

% Anonymised in this standalone copy. No section in the TMLR build
% uses either macro, so nothing renders from them; they stay defined
% so reinstating a section cannot break the build, and stay anonymous
% so the uploaded source leaks no identity. Restore the real values
% from paper/preamble.tex at camera-ready.
\newcommand{\repohost}{anonymous.repository.invalid}
\newcommand{\repourl}{https://\repohost}

% Shorthand macros used throughout.
\newcommand{\R}{\mathbb{R}}
\newcommand{\E}{\mathbb{E}}
\newcommand{\deff}{d_{\mathrm{eff}}}
\newcommand{\tauH}{\bar{\tau}_H}
\DeclareMathOperator{\rank}{rank}
\DeclareMathOperator{\range}{range}
\DeclareMathOperator{\tr}{tr}

% Shared theorem counter: Lemmas 1-2, Theorems 3-4, Remarks 5-6.
% Prose references depend on this numbering; do not split the counters.
\newtheorem{theorem}{Theorem}
\newtheorem{assumption}[theorem]{Assumption}
\newtheorem{lemma}[theorem]{Lemma}
\newtheorem{corollary}[theorem]{Corollary}
\newtheorem{claim}[theorem]{Claim}
\newtheorem{proposition}[theorem]{Proposition}
\theoremstyle{definition}
\newtheorem{definition}[theorem]{Definition}
\newtheorem{problem}{Problem}
\theoremstyle{remark}
\newtheorem{remark}[theorem]{Remark}
